\begin{abstract}
We study finite-horizon Markov Decision Processes with continuous states, finite actions, and a static Average Value-at-Risk (AVaR) objective. The augmented-state formulation enables dynamic programming but introduces an additional continuous state variable, making exact solutions computationally expensive. We propose a probabilistic Approximate Dynamic Programming (ADP) method using Gaussian processes to approximate Q-functions in the difficult region of the augmented state space, while exploiting risk-neutral continuation elsewhere. Bellman backups are allocated using acquisition scores based on decision uncertainty, reachability, and tail relevance. We also derive approximation and policy-performance error bounds. Numerical experiments show that decision-aware sampling often achieves better policy performance per Bellman backup than random or uncertainty-based sampling. However, incorporating reachability or tail relevance does not consistently improve performance, and greater backup efficiency does not necessarily imply lower runtime. These results suggest that decision-aware Bellman allocation is a promising approach to computationally efficient risk-averse control.
\end{abstract}